Боузер~--- злодей, который часто похищает принцесс, и даже имеет для их заточения целое подземелье. 

В подземелье расположено $n$ тюремных камер. Камеры пронумерованы натуральными числами от $1$ до $n$ и расположены в этом порядке вдоль длинного коридора на одинаковом расстоянии друг от друга. Камера с номером $i$ рассчитана на заточение $a_i$ принцесс. Вход в подземелье расположен между камерами $k$ и $k + 1$, при этом, если $k = 0$, это означает, что вход находится перед первой камерой. Если $k = n$, вход расположен после последней камеры.

Боузер собирается сделать $m$ вылазок с целью похищения принцесс. В $j$-ю вылазку он планирует похитить $b_j$ принцесс. Приведя их в подземелье, 
Боузер выбирает камеру для новых принцесс по следующему принципу:
\begin{itemize}
\item Если есть свободная камера с вместимостью $a_i = b_j$, то он обязательно выбирает такую камеру. Иначе он выбирает свободную камеру с вместимостью $a_i > b_j$.
\item Из подходящих камер он выбирает ту, которая расположена как можно ближе ко входу. Если камер с одинаковым расстоянием до входа несколько, он выбирает ту, у которой номер меньше. Расстояние от входа до камеры Боузер считает равным количеству камер между ними.
\item Если свободной камеры, подходящей по вместимости, нет, то ему остается лишь расстроиться и отпустить принцесс.
\end{itemize}
Чтобы не допустить побега принцесс, после того как Боузер заточил их в камеру, он закрывает камеру на ключ и выбрасывает его. Таким образом открыть камеру и добавить туда принцесс похищенных в другой вылазке, невозможно. Также Боузер никогда не пытается разместить принцесс, похищенных в одной вылазке, более чем в одной камере.

Марио и Луиджи нашли планы подземелья и владеют информацией о вылазках Боузера. Для каждой вылазки, они хотят определить камеру, в которую будут заточены принцессы.